\chapter{把基础用于后续课程 review}

\section{统一的检查顺序}
先确定对象和环境，再展开关键量词，检查存在唯一性，然后检查每一步推导及公理依赖。
完成后，把仍未证明的命题、例子提供的证据和已经证明的定理分开记录。

\begin{description}
\item[单变量与多元分析] 区分逐点与一致的量词顺序，写出确界存在的非空、有界条件，
  在构造极限对象时说明它属于哪个集合或空间。
\item[线性代数与代数] 写出底域、向量空间及映射的定义域和陪域；
  商上的运算要证明良定义，无限维存在性证明要标出选择依赖。
\item[测度与概率] 区分点、事件和事件族；\(\sigma\)-代数是 \(\mathcal P(X)\) 的子集，
  其成员才是 \(X\) 的子集。随机变量是指定可测空间之间的函数。
\item[拓扑] 明确环境集合、拓扑及子空间结构；区分「某个指定开覆盖」与「任意开覆盖」。
\item[计算机课程] 区分语法表达式、执行语义和关于程序的断言；
  用循环不变量说明归纳步骤，用终止条件区分部分正确性和全正确性。
\end{description}

\section{第一轮练习清单}
\begin{enumerate}
\item 写出 \(\forall x\exists y\,R(x,y)\) 的否定，并构造它与交换量词后的公式真值不同的结构。
\item 对一个含自由变量的公式做无捕获代入；解释错误代入改变了什么。
\item 从分离与外延证明交集存在且唯一，明确空索引族所用的环境集合。
\item 在 Kuratowski 编码中验证 \((a,a)=(b,c)\) 的情形。
\item 用最小归纳集合证明归纳原理，再给出一个有效递归定义。
\item 检查一个商集合上的运算是否与代表元选择无关。
\item 重写 Cantor 对角证明，标出分离的使用位置。
\item 从已有课程中选一个证明，列出假设、规则、所用定理和选择依赖。
\end{enumerate}

每次实际 review 以一小节或一个命题为单位；这些基础概念在真实证明中的用法，
会反过来决定哪些原书证明需要进一步展开。
